% 单位阶跃函数（综述）
% license CCBYSA3
% type Wiki

本文根据 CC-BY-SA 协议转载翻译自维基百科\href{https://en.wikipedia.org/wiki/Heaviside_step_function}{相关文章}

海维赛德阶跃函数，又称单位阶跃函数，通常记作$H$或$\theta$有时也记作 $u$; 1或$\mathbb{1}$是一种以英国工程师奥利弗·海维赛德命名的阶跃函数。它在自变量为负时取值为0在自变量为正时取值为1。关于函数在零点处的取值$H(0)$存在不同的约定。该函数是阶跃函数类的一种典型代表，所有的阶跃函数都可以表示为它的平移的线性组合。

这一函数最初出现在运算微积分中，用于求解微分方程。在该背景下，它表示一个在某一时刻“打开”并无限期保持“打开”的信号。海维赛德在分析电报通信时发展了运算微积分，并将该函数表示为1。  
\subsection{表述}
取约定 $H(0) = 1$，海维赛德函数可以定义为：

\begin{itemize}
\item 分段函数：
$$
H(x) := 
\begin{cases} 
1, & x \geq 0 \\[6pt] 
0, & x < 0 
\end{cases}~
$$
\item 使用艾弗森括号记号：
$$
H(x) := [x \geq 0]~
$$
\item 作为示性函数：
$$
H(x) := \mathbf{1}_{x \geq 0} = \mathbf{1}_{\mathbb{R}_{+}}(x)~
$$
若取另一种约定 $H(0) = \tfrac{1}{2}$，则可以表示为：
\item 分段函数：
$$
H(x) := 
\begin{cases} 
1, & x > 0 \\[6pt] 
\tfrac{1}{2}, & x = 0 \\[6pt] 
0, & x < 0 
\end{cases}~
$$
\item 符号函数的线性变换：
$$
H(x) := \tfrac{1}{2}\left(\operatorname{sgn}(x) + 1\right)~
$$
\item 两个艾弗森括号的算术平均：
$$
H(x) := \frac{[x \geq 0] + [x > 0]}{2}~
$$
双参数反正切函数（atan2）的单侧极限：

$$
H(x) := \lim_{\epsilon \to 0^{+}} \frac{\operatorname{atan2}(\epsilon, -x)}{\pi}
$$

**作为超函数（hyperfunction）：**

$$
H(x) := \left( 1 - \frac{1}{2\pi i}\log z, \; -\frac{1}{2\pi i}\log z \right)
$$

或等价地：

$$
H(x) := \left( -\frac{\log(-z)}{2\pi i}, \; -\frac{\log(-z)}{2\pi i} \right),
$$

其中 $\log z$ 表示复对数的主值。

---

其它在 $H(0)$ 处未定义的形式包括：

**分段函数：**

$$
H(x) := 
\begin{cases} 
1, & x > 0 \\[6pt] 
0, & x < 0 
\end{cases}
$$

**斜坡函数（ramp function）的导数：**

$$
H(x) := \frac{d}{dx}\max\{x,0\}, \quad x \neq 0
$$

**用绝对值函数表示：**

$$
H(x) = \frac{x + |x|}{2x}
$$

---

要不要我顺便帮你整理成 **对比表格**（不同约定下的多种定义方式并排对照），这样会更清晰直观？

\end{itemize}